% GENERATED FILE. Edit canonical lessons, then run npm run generate:books.
% canonical-source-hash: fnv1a64:65ea13036db166e1
% canonical-lessons: KA-C443-sancara, KA-C443-signal, KA-C443-padacari, KA-C443-dattane, KA-C443-niyantrisu

\chapter{City Traffic}
\label{ch:ka-city-traffic}
\hlchaptermodality{443}

\begin{chapteropening}
\textbf{By the end of this chapter:} \emph{I can say traffic, a traffic signal, a pedestrian, congestion and to control.}

Complete the last lesson of chapter 443: I can say traffic, a traffic signal, a pedestrian, congestion and to control.
\end{chapteropening}

\section[\texorpdfstring{\kn{ಸಂಚಾರ} (sañcāra)}{ಸಂಚಾರ (sañcāra)}]{\kn{ಸಂಚಾರ} (sañcāra) — traffic; movement}

\label{lesson:KA-C443-sancara}



\begin{quote}
Before the new one: say the Kannada for a mountain pass, then the Kannada for a route.
\end{quote}

\subsection*{You'll want to know: \kn{ಸಂಚಾರ}}
\textbf{\kn{ಸಂಚಾರ}} — \emph{sañcāra} — \textquotedblleft{}traffic; movement\textquotedblright{}.

\begin{quote}
\kn{ಸಂಜೆ} \kn{ಆರು} \kn{ಗಂಟೆಗೆ} \kn{ಸಂಚಾರ} \kn{ತುಂಬಾ} \kn{ನಿಧಾನ}.
\end{quote}

\begin{quote}
\emph{sañje āru gaṇṭege sañcāra tumbā nidhāna.}
\end{quote}

\begin{quote}
At six in the evening the traffic is very slow.
\end{quote}

\kn{ಸಂಚಾರ} \kn{ಪೊಲೀಸ್}, the traffic police; \kn{ಸಂಚಾರ} \kn{ನಿಯಮ}, a traffic rule.

\begin{cousinweb}
\textbf{\kn{ಸಂಚಾರ}} (\emph{saṃcāra}) is borrowed from Sanskrit \emph{saṃcāra}, \textquotedblleft{}walking about, a passage\textquotedblright{} (CDIAL 12869), from \emph{sam-} \textquotedblleft{}together\textquotedblright{} and \emph{car} \textquotedblleft{}to move\textquotedblright{}.
\end{cousinweb}

\begin{grammarlens}[title={What you've built}]
One more word for getting around.
\end{grammarlens}

\subsection*{Your turn}
\begin{itemize}
\raggedright
  \item \emph{Say it:} \emph{sañcāra}
  \item \emph{Say it:} \emph{sañcāra}, once more
  \item \emph{Say it:} \emph{sañje āru gaṇṭege sañcāra tumbā nidhāna.} — At six in the evening the traffic is very slow.
  \item \emph{Recall:} say the Kannada for a mountain pass, then the Kannada for a route, then say \emph{sañcāra} again
\end{itemize}

\begin{tcolorbox}[breakable,colback=teal!4,colframe=teal!35!black,title={Before you move on}]
What does \kn{ಸಂಚಾರ} mean? (\textquotedblleft{}Traffic; movement\textquotedblright{}.) Say it once more.
\end{tcolorbox}
\section[\texorpdfstring{\kn{ಸಿಗ್ನಲ್} (signal)}{ಸಿಗ್ನಲ್ (signal)}]{\kn{ಸಿಗ್ನಲ್} (signal) — a traffic signal}

\label{lesson:KA-C443-signal}



\begin{quote}
Before the new one: say the Kannada for a route, then the Kannada for traffic.
\end{quote}

\subsection*{You'll want to know: \kn{ಸಿಗ್ನಲ್}}
\textbf{\kn{ಸಿಗ್ನಲ್}} — \emph{signal} — \textquotedblleft{}a traffic signal\textquotedblright{}.

\begin{quote}
\kn{ಕೆಂಪು} \kn{ಸಿಗ್ನಲ್ನಲ್ಲಿ} \kn{ವಾಹನಗಳು} \kn{ನಿಲ್ಲಬೇಕು}.
\end{quote}

\begin{quote}
\emph{kempu signalnalli vāhanagaḷu nillabēku.}
\end{quote}

\begin{quote}
Vehicles must stop at a red signal.
\end{quote}

\kn{ಹಸಿರು} \kn{ಸಿಗ್ನಲ್}, a green light; \kn{ಸಿಗ್ನಲ್} \kn{ಸಿಗುವುದಿಲ್ಲ} is also said of a phone with no signal.

\begin{cousinweb}
\textbf{\kn{ಸಿಗ್ನಲ್}} (\emph{signal}) is a loan from English \emph{signal}.
\end{cousinweb}

\begin{grammarlens}[title={What you've built}]
One more word for getting around.
\end{grammarlens}

\subsection*{Your turn}
\begin{itemize}
\raggedright
  \item \emph{Say it:} \emph{signal}
  \item \emph{Say it:} \emph{signal}, once more
  \item \emph{Say it:} \emph{kempu signalnalli vāhanagaḷu nillabēku.} — Vehicles must stop at a red signal.
  \item \emph{Recall:} say the Kannada for a route, then the Kannada for traffic, then say \emph{signal} again
\end{itemize}

\begin{tcolorbox}[breakable,colback=teal!4,colframe=teal!35!black,title={Before you move on}]
What does \kn{ಸಿಗ್ನಲ್} mean? (\textquotedblleft{}A traffic signal\textquotedblright{}.) Say it once more.
\end{tcolorbox}
\section[\texorpdfstring{\kn{ಪಾದಚಾರಿ} (pādacāri)}{ಪಾದಚಾರಿ (pādacāri)}]{\kn{ಪಾದಚಾರಿ} (pādacāri) — a pedestrian}

\label{lesson:KA-C443-padacari}



\begin{quote}
Before the new one: say the Kannada for traffic, then the Kannada for a traffic signal.
\end{quote}

\subsection*{You'll want to know: \kn{ಪಾದಚಾರಿ}}
\textbf{\kn{ಪಾದಚಾರಿ}} — \emph{pādacāri} — \textquotedblleft{}a pedestrian\textquotedblright{}.

\begin{quote}
\kn{ಪಾದಚಾರಿಗಳು} \kn{ರಸ್ತೆಯ} \kn{ಅಂಚಿನಲ್ಲಿ} \kn{ನಡೆಯಬೇಕು}.
\end{quote}

\begin{quote}
\emph{pādacārigaḷu rasteya añcinalli naḍeyabēku.}
\end{quote}

\begin{quote}
Pedestrians should walk at the edge of the road.
\end{quote}

Plural: \kn{ಪಾದಚಾರಿಗಳು}, \emph{pādacārigaḷu}.

\kn{ಪಾದ}, a foot, plus -\kn{ಚಾರಿ}, one who goes: one who goes on foot.

\begin{cousinweb}
\textbf{\kn{ಪಾದಚಾರಿ}} (\emph{pādacāri}) is borrowed from Sanskrit \emph{pādacārin}, \textquotedblleft{}going on foot, a pedestrian\textquotedblright{} (Monier-Williams), from \emph{pāda}, \textquotedblleft{}a foot\textquotedblright{}, and \emph{cārin}, \textquotedblleft{}one who moves\textquotedblright{}.
\end{cousinweb}

\begin{grammarlens}[title={What you've built}]
One more word for getting around.
\end{grammarlens}

\subsection*{Your turn}
\begin{itemize}
\raggedright
  \item \emph{Say it:} \emph{pādacāri}
  \item \emph{Say it:} \emph{pādacāri}, once more
  \item \emph{Say it:} \emph{pādacārigaḷu rasteya añcinalli naḍeyabēku.} — Pedestrians should walk at the edge of the road.
  \item \emph{Recall:} say the Kannada for traffic, then the Kannada for a traffic signal, then say \emph{pādacāri} again
\end{itemize}

\begin{tcolorbox}[breakable,colback=teal!4,colframe=teal!35!black,title={Before you move on}]
What does \kn{ಪಾದಚಾರಿ} mean? (\textquotedblleft{}A pedestrian\textquotedblright{}.) Say it once more.
\end{tcolorbox}
\section[\texorpdfstring{\kn{ದಟ್ಟಣೆ} (daṭṭaṇe)}{ದಟ್ಟಣೆ (daṭṭaṇe)}]{\kn{ದಟ್ಟಣೆ} (daṭṭaṇe) — congestion; crowding}

\label{lesson:KA-C443-dattane}



\begin{quote}
Before the new one: say the Kannada for a traffic signal, then the Kannada for a pedestrian.
\end{quote}

\subsection*{You'll want to know: \kn{ದಟ್ಟಣೆ}}
\textbf{\kn{ದಟ್ಟಣೆ}} — \emph{daṭṭaṇe} — \textquotedblleft{}congestion; crowding\textquotedblright{}.

\begin{quote}
\kn{ರಸ್ತೆಯಲ್ಲಿ} \kn{ವಾಹನ} \kn{ದಟ್ಟಣೆ} \kn{ಇತ್ತು}.
\end{quote}

\begin{quote}
\emph{rasteyalli vāhana daṭṭaṇe ittu.}
\end{quote}

\begin{quote}
There was a traffic jam on the road.
\end{quote}

\kn{ವಾಹನ} \kn{ದಟ್ಟಣೆ}, a traffic jam; \kn{ದಟ್ಟ} means dense or thick.

\begin{cousinweb}
\textbf{\kn{ದಟ್ಟಣೆ}} (\emph{daṭṭaṇe}) is native Dravidian, built on Kannada \emph{daṭṭa}, \textquotedblleft{}thick, crowded, dense\textquotedblright{} (DEDR 3020).
\end{cousinweb}

\begin{grammarlens}[title={What you've built}]
One more word for getting around.
\end{grammarlens}

\subsection*{Your turn}
\begin{itemize}
\raggedright
  \item \emph{Say it:} \emph{daṭṭaṇe}
  \item \emph{Say it:} \emph{daṭṭaṇe}, once more
  \item \emph{Say it:} \emph{rasteyalli vāhana daṭṭaṇe ittu.} — There was a traffic jam on the road.
  \item \emph{Recall:} say the Kannada for a traffic signal, then the Kannada for a pedestrian, then say \emph{daṭṭaṇe} again
\end{itemize}

\begin{tcolorbox}[breakable,colback=teal!4,colframe=teal!35!black,title={Before you move on}]
What does \kn{ದಟ್ಟಣೆ} mean? (\textquotedblleft{}Congestion; crowding\textquotedblright{}.) Say it once more.
\end{tcolorbox}
\section[\texorpdfstring{\kn{ನಿಯಂತ್ರಿಸು} (niyantrisu)}{ನಿಯಂತ್ರಿಸು (niyantrisu)}]{\kn{ನಿಯಂತ್ರಿಸು} (niyantrisu) — to control; to regulate}

\label{lesson:KA-C443-niyantrisu}



\begin{quote}
Before the new one: say the Kannada for a pedestrian, then the Kannada for congestion.
\end{quote}

\subsection*{You'll want to know: \kn{ನಿಯಂತ್ರಿಸು}}
\textbf{\kn{ನಿಯಂತ್ರಿಸು}} — \emph{niyantrisu} — \textquotedblleft{}to control; to regulate\textquotedblright{}.

\begin{quote}
\kn{ಒಬ್ಬ} \kn{ಪೊಲೀಸ್} \kn{ಸಂಚಾರವನ್ನು} \kn{ನಿಯಂತ್ರಿಸುತ್ತಿದ್ದನು}.
\end{quote}

\begin{quote}
\emph{obba polīs sañcāravannu niyantrisuttiddanu.}
\end{quote}

\begin{quote}
A policeman was controlling the traffic.
\end{quote}

Past (I …): \kn{ನಿಯಂತ್ರಿಸಿದೆ}, \emph{niyantriside}.

What is controlled takes -\kn{ಅನ್ನು}; \kn{ನಿಯಂತ್ರಣ} is control.

\begin{cousinweb}
\textbf{\kn{ನಿಯಂತ್ರಿಸು}} (\emph{niyantrisu}) is built from Sanskrit \emph{niyantraṇa}, \textquotedblleft{}restraining, checking, governing\textquotedblright{} (Monier-Williams), with the verb ending -\emph{isu}.
\end{cousinweb}

\begin{grammarlens}[title={What you've built}]
One more word for getting around.
\end{grammarlens}

\subsection*{Your turn}
\begin{itemize}
\raggedright
  \item \emph{Say it:} \emph{niyantrisu}
  \item \emph{Say it:} \emph{niyantrisu}, once more
  \item \emph{Say it:} \emph{obba polīs sañcāravannu niyantrisuttiddanu.} — A policeman was controlling the traffic.
  \item \emph{Recall:} say the Kannada for a pedestrian, then the Kannada for congestion, then say \emph{niyantrisu} again
\end{itemize}

\begin{tcolorbox}[breakable,colback=teal!4,colframe=teal!35!black,title={Before you move on}]
What does \kn{ನಿಯಂತ್ರಿಸು} mean? (\textquotedblleft{}To control; to regulate\textquotedblright{}.) Say it once more.
\end{tcolorbox}
